\begin{table}[H]
\centering
\caption{Effect of Reform Exposure on Labor Market Outcomes}
\begin{threeparttable}
\begin{tabular}{lccc}
\toprule
 & (1) & (2) & (3) \\
 & Temp.\ Share & Log Emp. & Temp.\ Share \\
 & (Unweighted) & (Unweighted) & (Pop.-Weighted) \\
\midrule
Pre-Reform Temp Share $\times$ Post & -0.2200 & 0.1948 & -0.4622$^{***}$ \\
  & (0.2049) & (0.2962) & (0.0460) \\
Wild bootstrap $p$-value & [0.362] & [0.564] & [0.009] \\
\midrule
Region FE & \checkmark & \checkmark & \checkmark \\
Quarter FE & \checkmark & \checkmark & \checkmark \\
Population weights & & & \checkmark \\
Observations & 1,140 & 1,140 & 1,140 \\
Regions & 19 & 19 & 19 \\
Dep.\ Var.\ Mean & 0.236 & 6.103 & 0.236 \\
\bottomrule
\end{tabular}
\begin{tablenotes}[flushleft]
\small
\item Notes: Each column reports the coefficient on the interaction of pre-reform regional temporary employment share (2021 average) with a post-reform indicator (2022Q2 onwards). Standard errors clustered at the region level in parentheses. Wild cluster bootstrap $p$-values (Webb weights, 9,999 iterations) in brackets. Column~(3) uses 2021 average total wage earners as population weights. Since temporary and permanent employment shares sum to unity by construction, the permanent share coefficient is mechanically equal in magnitude and opposite in sign to Column~(1); we omit it for parsimony. $^{***}$ $p<0.01$, $^{**}$ $p<0.05$, $^{*}$ $p<0.10$.
\end{tablenotes}
\end{threeparttable}
\label{tab:main}
\end{table}

